% Section: Abstract (body only; the abstract environment lives in main.tex)
The partition function $Z$ defines the thermal equilibrium states of atomic systems, yet its evaluation for high-dimensional heteroepitaxial growth remains a persistent analytical hurdle. In the $\text{Fe/MgO(001)}$ system, this difficulty is compounded by surface free energy mismatches that drive the island (Volmer-Weber) growth mode rather than the preferred flat (Frank–van der Merwe) mode. To address this, we introduce a biased exploration within an active learning framework to sample an ensemble of surrogate-relaxed structures. Within this ensemble, the island mode is identified as the ground state, consistent with experimental observations. Applying dimensionality reduction and Kernel Density Estimation (KDE) to this ensemble, we categorize two distinct growth modes: island and flat. Using electronic structure analysis, we demonstrate that the island mode is preferred due to localized interfacial $\text{Fe}$ $3d$ orbitals near the Fermi level, which destabilize the flat mode. Furthermore, we determine that while higher temperatures raise the probability of the flat mode according to Boltzmann statistics, the energy barrier prevents thermal fluctuations from suppressing the island mode, necessitating non-thermal strategies to engineer the flat interface.